\subsection*{5.5.}
上述纤维化为曲线的方法还给出 Frobenius 特征值的复绝对值的 Lang--Weil 估计~[3].

设 $X$ 是 $\mathbb F_q$ 上的有限型概形, $\underline G$ 是 $X$ 上的可构造 $E$ 层,
$\lambda\in\mathbb R$. 假设对每个点 $x\in X$, 令 $\#k(x)=q_x$,
$F_x$ 在 $G_x$ 上的特征值均为代数数 $\alpha$, 其所有复绝对值均满足
\[
|\alpha|\leq q_x^\lambda.
\]
此时记作
\[
\log|\underline G|\leq\lambda.
\]

\subsubsection*{引理 5.5.1.}
设 $X$ 是 $\mathbb F_q$ 上的曲线, 此处指 $\mathbb F_q$ 上维数至多为 $1$
的有限型分离概形 $X$. 设 $\underline G$ 是 $X$ 上的可构造 $E$ 层,
满足 $\log|\underline G|\leq\lambda$, 其中 $\lambda\in\mathbb R$.
则 $F$ 在 $\underline H_c^1(X,G)$ 上的特征值 $\alpha$ 的复绝对值满足
\[
|\alpha|\leq q^{\lambda+1}.
\]

\medskip
\noindent\emph{证明.}
设 $\sigma$ 是 $E$ 的复嵌入, 或仅仅是 $E$ 中 $\mathbb Q$ 的代数闭包的复嵌入.
只需验证对每个 $\sigma$, 有理函数 $\sigma Z(X,\underline G,t)$ 的零点 $\beta$ 满足
\[
|\beta^{-1}|\leq q^{\lambda+1}.
\]
事实上, 引理~5.2.1 已表明, 若 $\alpha$ 是 Frobenius 在
$\underline H_c^0(X,\underline G)$ 或 $\underline H_c^2(X,\underline G)$ 上的特征值,
则 $|\sigma\alpha|\leq q^{\lambda+1}$.

我们验证级数
\[
\log\sigma Z(X,\underline G,t)
=\sum_k\frac{t^k}{k}
  \sum_{x\in X(\mathbb F_{q^k})}
  \sigma\operatorname{Tr}(F,x^*\underline G)
\]
在 $|t|<q^{-\lambda-1}$ 时收敛. 将 $X$ 表示为 $\mathbb P^1$ 的分歧覆盖,
可得初等估计
\[
\#X(\mathbb F_{q^k})\leq A q^k.
\]
另一方面, 假设给出
\[
\left|\sigma\operatorname{Tr}(F,x^*\underline G)\right|
\leq B(q^k)^\lambda.
\]
因而 $\log\sigma Z(X,\underline G,t)$ 中 $t^k$ 的系数 $c_k$ 满足
\[
|c_k|\leq\frac{1}{k}\,C(q^{1+\lambda})^k,
\]
断言随之得出.

\subsubsection*{定理 5.5.2.}
设 $X$ 是 $\mathbb F_q$ 上维数至多为 $n$ 的有限型分离概形,
$\underline G$ 是 $X$ 上的可构造 $E$ 层, 满足 $\log|\underline G|\leq\lambda$,
其中 $\lambda\in\mathbb R$. 则 Frobenius 在 $\underline H_c^i(X,\underline G)$ 上的
特征值 $\alpha$ 的复绝对值满足
\[
|\alpha|\leq q^{\lambda+i}
\qquad\text{且}\qquad
|\alpha|\leq q^{\lambda+n}.
\]

\medskip
\noindent\emph{证明.}
若 $n\leq0$, 定理是平凡的. 若 $n=1$, 其 $i=0$ 和 $i=2$ 的情形由引理~5.2.1 得出,
而 $i=1$ 的情形由引理~5.5.1 得出. 随后如定理~5.2.2 的证明那样对 $n$ 归纳.

\subsubsection*{推论 5.5.3.}
设 $X$ 是 $\mathbb F_q$ 上维数为 $n$ 的有限型分离概形.
则 $\underline H_c^i(X,\mathbb Q_\ell)$ 的几何 Frobenius 自同态 $F$
的特征值 $\alpha$ 均为代数整数. 此外:
\begin{enumerate}[label=(\roman*)]
\item 对 $\ell'\ne p$, $\alpha$ 是 $\ell'$ 进单位;
\item 对 $i>n$, $\alpha$ 被 $q^{i-n}$ 整除;
\item 对 $i\leq n$ (分别地, $i\geq n$),
  $q^i\alpha^{-1}$ (分别地, $q^n\alpha^{-1}$) 是代数整数;
\item 对 $0<i$, $\alpha$ 的复绝对值满足
  \[
  |\alpha|\leq q^{i-\frac12}
  \qquad\text{且}\qquad
  |\alpha|\leq q^n.
  \]
\end{enumerate}

\medskip
\noindent\emph{证明.}
由定理~5.2.2, $\alpha$ 和 $q^{n-i}\alpha$ 均为代数整数.
断言 (i)--(iii) 随即由定理~5.4 得出. 当 $i>n$ 时, 断言 (iv) 由定理~5.5.2 得出.
要在 $i\leq n$ 时得到 (iv), 需将定理~5.5.2 与曲线的 Riemann 假设结合.

对 $n$ 归纳, $n=0$ 的情形是平凡的. 若 $U$ 为 $X$ 的开子集,
其补集 $Z$ 维数小于 $n$, 则正合列
\[
\underline H_c^i(U,\mathbb Q_\ell)\longrightarrow
\underline H_c^i(X,\mathbb Q_\ell)\longrightarrow
\underline H_c^i(Z,\mathbb Q_\ell)
\]
表明只需对 $U$ 证明所求断言. 先将 $X$ 换成 $X_{\mathrm{red}}$ 后,
这使我们归结为 $X$ 光滑仿射, 且存在满足下述条件的光滑态射 $f:X\longrightarrow S$ 的情形:
\begin{enumerate}[label=(\alph*)]
\item $\dim S\leq n-1$;
\item $f$ 的纤维为纯一维的.
\end{enumerate}
由于 $f$ 的纤维为仿射曲线, 有
\[
R^0f_!\mathbb Q_\ell=0.
\]
由 Weil~[4],
\begin{equation}\tag{5.5.3.1}
\log\bigl|R^1f_!\mathbb Q_\ell\bigr|\leq\frac12,
\end{equation}
而将引理~5.2.1 用于 $f$ 的纤维得到
\begin{equation}\tag{5.5.3.2}
\log\bigl|R^2f_!\mathbb Q_\ell\bigr|\leq1.
\end{equation}
$f$ 的 Leray 谱序列在此退化为长正合列; 相关部分为
\[
\underline H_c^{i-2}(S,R^2f_!\mathbb Q_\ell)\longrightarrow
\underline H_c^i(X,\mathbb Q_\ell)\longrightarrow
\underline H_c^{i-1}(S,R^1f_!\mathbb Q_\ell).
\]
由 (5.5.3.1) (分别地, (5.5.3.2)) 及定理~5.5.2, Frobenius 在
$\underline H_c^{i-1}(S,R^1f_!\mathbb Q_\ell)$
(分别地, $\underline H_c^{i-2}(S,R^2f_!\mathbb Q_\ell)$) 上的特征值 $\alpha$
的复绝对值满足
\[
|\alpha|\leq q^{i-1+\frac12}=q^{i-\frac12}
\]
(分别地,
\[
|\alpha|\leq q^{i-2+1}=q^{i-1}\leq q^{i-\frac12}).
\]
推论得证.

\medskip
\noindent\emph{底本注.}
照录法文层保留了三处印误: 将系数 $c_k$ 说成属于 $\sigma Z$,
但所列级数及因子 $1/k$ 表明它是 $\log\sigma Z$ 的系数;
定理 5.5.2 的证明引用自身, 而非定理 5.2.2 中的平行归纳;
推论 5.5.3 的证明说要对 $U$ 证明定理 5.2.2, 但该论证实际证明的是推论的断言 (iv).
校订法文层和英文层给出相容的读法. 直接对照扫描件也补回了定理 5.5.2 的陈述中
继承的法文 TeX 所漏的记号下划线.
